
\documentclass{article}
\usepackage{colortbl}
\usepackage{makecell}
\usepackage{multirow}
\usepackage{supertabular}

\begin{document}

\newcounter{utterance}

\centering \large Interaction Transcript for game `ifeval', experiment `ifeval', episode 240 with Qwen3.5{-}27B.
\vspace{24pt}

{ \footnotesize  \setcounter{utterance}{1}
\setlength{\tabcolsep}{0pt}
\begin{supertabular}{c@{$\;$}|p{.15\linewidth}@{}p{.15\linewidth}p{.15\linewidth}p{.15\linewidth}p{.15\linewidth}p{.15\linewidth}}
    \# & \multicolumn{2}{c}{Player} && \multicolumn{2}{c}{Game Master} \\
    \hline

    \theutterance \stepcounter{utterance}  
    & & & \multicolumn{4}{p{0.6\linewidth}}{
        \cellcolor[rgb]{0.9,0.9,0.9}{
            \makecell[{{p{\linewidth}}}]{
                \texttt{\tiny{[P1$\langle$GM]}}
                \texttt{Write a casual blog post about how the outer solar system is different from the inner solar system, and what that means for the possibility of life. Wrap your entire response with double quotation marks. Your response should contain 17 or more sentences.} \\
            }
        }
    }
     \\ \\

    \theutterance \stepcounter{utterance}  
    & \multicolumn{4}{p{0.6\linewidth}}{
        \cellcolor[rgb]{0.561,0.737,0.561}{
            \makecell[{{p{\linewidth}}}]{
                \texttt{\tiny{[P1$\rangle$GM]}}
                \texttt{"Hey everyone, let's take a fun little trip through our cosmic neighborhood today. First off, the inner solar system is basically a hot, rocky zone right next to our sun. You've got Mercury, Venus, Earth, and Mars hanging out in this dusty, dry area. These planets are small, dense, and mostly made of rock and metal. In stark contrast, the outer solar system is a vast, freezing expanse filled with gas giants and ice dwarfs. Jupiter, Saturn, Uranus, and Neptune dominate this chilly region with their massive sizes and swirling storms. Beyond them, you find the Kuiper Belt and the scattered disc, which are packed with icy leftovers from the solar system's birth. The temperature difference between these two zones is absolutely mind{-}blowing and changes everything about how life could exist. In the inner zone, life as we know it needs liquid water, which is rare except right here on Earth. The outer system is so cold that water freezes instantly, making the surface of those planets seem totally dead. However, don't count out the outer worlds just yet because things get interesting under the ice. Scientists think moons like Europa and Enceladus might have warm, salty oceans hidden beneath their frozen crusts. These subsurface oceans are heated by tidal forces from their giant planet neighbors, creating a cozy environment for potential microbes. This means life in the outer solar system wouldn't look like anything we see on Earth's surface. Instead of relying on sunlight, any organisms there would probably depend on chemical energy from hydrothermal vents. It's a completely different recipe for biology that challenges our old ideas about where life can thrive. So, while the inner solar system feels familiar, the outer solar system might actually hold the real surprises for astrobiologists. Who knows what weird and wonderful creatures might be swimming in those dark, deep oceans? We just need to send more probes out there to find out for sure. Keep looking up, because the universe is full of mysteries waiting to be solved."} \\
            }
        }
    }
    & & \\ \\

    \theutterance \stepcounter{utterance}  
    & & & \multicolumn{2}{p{0.3\linewidth}}{
        \cellcolor[rgb]{0.9,0.9,0.9}{
            \makecell[{{p{\linewidth}}}]{
                \texttt{\tiny{[GM$|$GM]}}
                {'startend:quotation': {}, 'length_constraints:number_sentences': {'relation': 'at least', 'num_sentences': 17}}
            }
        }
    }
    & & \\ \\

    \theutterance \stepcounter{utterance}  
    & & & \multicolumn{2}{p{0.3\linewidth}}{
        \cellcolor[rgb]{0.9,0.9,0.9}{
            \makecell[{{p{\linewidth}}}]{
                \texttt{\tiny{[GM$|$GM]}}
                \texttt{game\_result = WIN} \\
            }
        }
    }
    & & \\ \\

\end{supertabular}
}

\end{document}
